% ----- CONCLUDING DOHA -----
\subsection*{\deva{अन्तिम दोहा} (Concluding Doha)}
\addcontentsline{toc}{subsection}{\deva{अन्तिम दोहा} (Concluding Doha) — \deva{पवन तनय...}}

\begin{minipage}[t]{0.55\textwidth}
\vspace{0pt}

{\large \linespread{1.5}\selectfont
\makebox[\linewidth][l]{\deva{पवन तनय संकट हरन}}\\
\makebox[\linewidth][c]{\deva{मंगल मूरति रूप।}}\\
\makebox[\linewidth][l]{\deva{राम लखन सीता सहित }}\\
\makebox[\linewidth][c]{\deva{हृदय बसहु सुर भूप॥}}
\par}

\vspace{0.5em}

{\large \linespread{1.5}\selectfont
\makebox[\linewidth][l]{\telu{పవన తనయ సంకట హరన} }\\
\makebox[\linewidth][c]{\telu{మంగల మూరతి రూప।}}\\
\makebox[\linewidth][l]{\telu{రామ లఖన సీతా సహిత}}\\
\makebox[\linewidth][c]{\telu{హృదయ బసహు సుర భూప॥}}
\par}

\vspace{0.5em}

{\large \linespread{1.3}\selectfont
\textit{pavana tanaya saṅkaṭa harana, maṅgala mūrati rūpa |}\\
\textit{rāma lakhana sītā sahita, hṛdaya basahu sura bhūpa ||}
\par}
\end{minipage}%
\hfill
\begin{minipage}[t]{0.42\textwidth}
\vspace{0pt}
\begin{tcolorbox}[anchorbox]
\textbf{The Crisis Preventer (Strategic Leadership):} Before returning to Ayodhya, Lord Rama sent Hanuman ahead to meet Bharata and gauge his true feelings. If Bharata had grown accustomed to power, Rama's sudden return could cause a political crisis or emotional distress. Hanuman acted with supreme diplomacy, emotional intelligence, and leadership, ensuring a seamless, joyful transition of power and preventing any potential crisis.
\vspace{0.5em}\newline Hanuman's strategic foresight and crisis leadership resolved a highly delicate situation before it could escalate.\newline\null\hfill\href{https://www.valmiki.iitk.ac.in/sloka?field_kanda_tid=6&language=dv&field_sarga_value=125}{\textit{Yuddha Kanda, 125}}
\end{tcolorbox}
\end{minipage}

\vspace{0.5em}
\subsubsection*{\textbf{Visual Rhythm Map:}}
\begin{table}[h!]
\centering
\renewcommand{\arraystretch}{1.2}
\begin{tabularx}{\textwidth}{|>{\centering\arraybackslash\hsize=0.8750\hsize}X|>{\centering\arraybackslash\hsize=0.8750\hsize}X|>{\centering\arraybackslash\hsize=1.1667\hsize}X|>{\centering\arraybackslash\hsize=0.8750\hsize}X|>{\centering\arraybackslash\hsize=1.1667\hsize}X|>{\centering\arraybackslash\hsize=1.1667\hsize}X|>{\centering\arraybackslash\hsize=0.8750\hsize}X|}
\hline
\textbf{\laghu{प}\laghu{व}\laghu{न}} \par (3) & \textbf{\laghu{त}\laghu{न}\laghu{य}} \par (3) & \textbf{\guru{सं}\laghu{क}\laghu{ट}} \par (4) & \textbf{\laghu{ह}\laghu{र}\laghu{न}} \par (3) & \textbf{\guru{मं}\laghu{ग}\laghu{ल}} \par (4) & \textbf{\guru{मू}\laghu{र}\laghu{ति}} \par (4) & \textbf{\guru{रू}\laghu{प}} \par (3) \\
\hline
\end{tabularx}
\vspace{-1.5em}
\end{table}

\begin{table}[h!]
\centering
\renewcommand{\arraystretch}{1.2}
\begin{tabularx}{\textwidth}{|>{\centering\arraybackslash\hsize=1.0000\hsize}X|>{\centering\arraybackslash\hsize=1.0000\hsize}X|>{\centering\arraybackslash\hsize=1.3333\hsize}X|>{\centering\arraybackslash\hsize=1.0000\hsize}X|>{\centering\arraybackslash\hsize=1.0000\hsize}X|>{\centering\arraybackslash\hsize=1.0000\hsize}X|>{\centering\arraybackslash\hsize=0.6667\hsize}X|>{\centering\arraybackslash\hsize=1.0000\hsize}X|}
\hline
\textbf{\guru{रा}\laghu{म}} \par (3) & \textbf{\laghu{ल}\laghu{ख}\laghu{न}} \par (3) & \textbf{\guru{सी}\guru{ता}} \par (4) & \textbf{\laghu{स}\laghu{हि}\laghu{त}} \par (3) & \textbf{\laghu{हृ}\laghu{द}\laghu{य}} \par (3) & \textbf{\laghu{ब}\laghu{स}\laghu{हु}} \par (3) & \textbf{\laghu{सु}\laghu{र}} \par (2) & \textbf{\guru{भू}\laghu{प}} \par (3) \\
\hline
\end{tabularx}
\vspace{-1.5em}
\end{table}

\vspace{1em}
\textbf{ \deva{सरल-संस्कृतम्}:}\\
 \deva{हे पवनपुत्र! हे सङ्कट-हारक! भवान् मङ्गल-मूर्तेः साक्षात् रूपम् अस्ति।}\\
 \deva{हे देवानां राजन्! भवान् राम-लक्ष्मण-सीताभिः सह सर्वदा मम हृदये वसतु॥}\\


O Son of the Wind, O Destroyer of Crisis! You are the very embodiment of auspiciousness.\\
O King of Gods! Please reside forever in my heart, along with Rama, Lakshmana, and Sita.


\vspace{1em}
\newpage
\textbf{\deva{प्रतिपदार्थः} (Meaning of each word):}
\begin{table}[h!]
\centering
\renewcommand{\arraystretch}{1.2}
\begin{tabularx}{\textwidth}{@{} l l X X @{}}
\toprule
\textbf{Awadhi} & \textbf{Sanskrit} & \textbf{English} & \textbf{Grammar \& Notes} \\
\midrule
\deva{पवन तनय} / \telu{పవన తనయ} / \textit{pavana tanaya}  &  \deva{पवन-तनय}  &  Son of the Wind-God  &   \\
\deva{संकट हरन} / \telu{సంకట హరన} / \textit{saṅkaṭa harana}  &  \deva{सङ्कट-हरण}  &  Destroyer of Crisis  &   \\
\deva{मंगल मूरति} / \telu{మంగల మూరతి} / \textit{maṅgala mūrati}  &  \deva{मङ्गल-मूर्तिः}  &  Auspicious idol/form  &   \\
\deva{रूप} / \telu{రూప} / \textit{rūpa}  &  \deva{रूपम्}  &  Embodiment  &   \\
\deva{राम लखन} / \telu{రామ లఖన} / \textit{rāma lakhana}  &  \deva{राम-लक्ष्मण}  &  Rama and Lakshmana  & \\ 
\deva{सीता सहित} / \telu{సీతా సహిత} / \textit{sītā sahita}  &  \deva{सीतया सह}  &  Along with Sita  &   \\
\deva{हृदय बसहु} / \telu{హృదయ బసహు} / \textit{hṛdaya basahu}  &  \deva{हृदये वसतु}  &  Reside in the heart  & Imperative `-hu` suffix \\
\deva{सुर भूप} / \telu{సుర భూప} / \textit{sura bhūpa}  &  \deva{देवानां राजा / भूपतिः}  &  King of Gods  &   \\
\bottomrule
\end{tabularx}
\end{table}

\vspace{0.5em}
\begin{tcolorbox}[poetrybox]
\textbf{Structural Symmetry:}\\
The Chalisa ends exactly how it began. 
By returning flawlessly to the beautiful 13-11 metric cadence of the \textbf{Doha} format, perfectly bookmarking the 40 marching Chaupais with reverence and praise.
\end{tcolorbox}

\newpage
